\documentclass{article}
\usepackage{amsmath, amssymb, amsthm}

\newtheorem*{theorem210}{Theorem 2.10}
\newtheorem*{theoremA7}{Theorem A.7}

\begin{document}

\begin{theorem210}
If the $n \times n$ matrix $A$ is strictly diagonally dominant, then:
\begin{enumerate}
    \item $A$ is a non-singular matrix.
    \item For every vector $b$ and every starting guess, the Jacobi method applied to $Ax = b$ converges to the (unique) solution.
\end{enumerate}
\end{theorem210}

\begin{proof}
Let $A = D + L + U$, where $D$ is the diagonal part of $A$, and $L$ and $U$ are the strictly lower and upper triangular parts, respectively. Let $R = L + U$ denote the non-diagonal part of the matrix $A$. Note that the diagonal entries of $R$ are zero.

To guarantee the convergence of the Jacobi method, we must show that the spectral radius of the iteration matrix is less than one, that, is $\rho(D^{-1}R) < 1$. In other words, we have to show that every eigenvalue of this matrix has magnitude strictly less than one.

Let $\lambda$ be an arbitrary eigenvalue of $D^{-1}R$ with corresponding eigenvector $v$. By definition, we have:
\[
D^{-1}Rv = \lambda v \implies Rv = \lambda Dv
\]

Choose the eigenvector $v$ such that its infinity norm is 1, i.e., $\|v\|_\infty = 1$. This means there exists an index $m$ ($1 \le m \le n$) such that the component $|v_m| = 1$, and for all other components, $|v_j| \le 1$.

Considering the $m$-th component of the vector equation $Rv = \lambda Dv$, and recalling that $r_{mm} = 0$ (because, as already stated, all diagonal values of $R$ are zero), we get:
\[
\sum_{\substack{j \neq m}}^{n} r_{mj}v_j = \lambda d_{mm}v_m
\]

Taking the absolute value of both sides and substituting $|v_m| = 1$, produces:
\[
|\lambda| |d_{mm}| = \left| \sum_{\substack{j \neq m}}^{n} r_{mj}v_j \right|\]

Applying the triangle inequality, we get:
\[
|\lambda| |d_{mm}| \le \sum_{\substack{j \neq m}}^{n} |r_{mj}| |v_j|
\]

Since $|v_j| \le 1$ for all $j$, we can bound the right side:
\[
|\lambda| |d_{mm}| \le \sum_{\substack{j \neq m}}^{n} |r_{mj}|
\]

By the hypothesis that $A$ is strictly diagonally dominant, the absolute value of the diagonal entry strictly dominates the sum of the absolute values of the off-diagonal entries in its row:
\[
\sum_{\substack{j \neq i}}^{n} |a_{ij}| < |a_{ii}| \implies \sum_{\substack{j \neq m}}^{n} |r_{mj}| < |d_{mm}|
\]

Combining these inequalities gives:
\[
|\lambda| |d_{mm}| < |d_{mm}|
\]

Since $A$ is strictly diagonally dominant, its diagonal elements cannot be zero, so $|d_{mm}| > 0$. Dividing both sides by $|d_{mm}|$ yields:
\[
|\lambda| < 1
\]

Because $\lambda$ was an arbitrary eigenvalue, it follows that $\rho(D^{-1}R) < 1$: this consequence is due to the fact that we have proven that all eigenvalues are less than one. We now rely on the general theorem on iterative methods:

\begin{theoremA7}
If the $n \times n$ matrix $M$ has spectral radius $\rho (M) < 1$, and $c$ is arbitrary, then, for any vector $x_0$, the iteration $x_{k+1} = Mx_k + c$ converges. In fact, there exists a unique $x_*$ such that $\lim_{k \to \infty} x_k = x_*$ and $x_* = Mx_* + c$.
\end{theoremA7}

Therefore, let $M = -D^{-1}(L + U)$ and $c = D^{-1}b$, we have:
\[
x_{k+1} = Mx_k + c = -D^{-1}(L + U)x_k + D^{-1}b
\]

This guarantees that the Jacobi method converges for any starting guess $x_0$ and any vector $b$.

Finally, since the linear system $Ax = b$ is guaranteed to have a unique solution for every possible $b$, fundamental linear algebra affirms that the matrix $A$ must be invertible. Therefore, $A$ is a non-singular matrix, completing the proof.
\end{proof}

\end{document}
